% ID: SEL-D0014-BD-E08
% ESTADO: validada. S015; tres pasadas del mismo autor, sin segundo revisor.
\textbf{a) Centro del intervalo dado.}\par\nopagebreak[4]
El intervalo normal de media con desviación conocida es simétrico alrededor de la media muestral. Por tanto,
\[
\boxed{\overline x=\frac{6.04+8}{2}=7.02\text{ cm}.}
\]
\textbf{b) Error del intervalo dado.}\par\nopagebreak[4]
Es su semiamplitud, no su ancho completo:
\[
\boxed{E_{36}=\frac{8-6.04}{2}=0.98\text{ cm}.}
\]
Coincide con $1.96\cdot3/\sqrt{36}=0.98$ cm usando el cuantil tabulado. Con el cuantil preciso se obtiene aproximadamente $0.979982$ cm, coherente con los extremos impresos redondeados.

\textbf{c) Tamaño 49 con la misma confianza 95\,\%.}\par\nopagebreak[4]
Al mantener confianza y desviación, $E=z\sigma/\sqrt n$. Aumentar $n$ reduce el error:
\[
E_{49}=1.96\frac3{\sqrt{49}}=\boxed{0.84\text{ cm}}.
\]
Con el cuantil preciso, $E_{49}\approx0.839984565$ cm. También se comprueba mediante $E_{49}/E_{36}=\sqrt{36/49}=6/7$. Esta reducción es de aproximadamente el 14.29\,\%. No hace falta suponer que la nueva muestra mantenga la media $7.02$ para calcular el error; no se construye un nuevo intervalo sin conocer su centro.

\textbf{d) Aumentar la confianza.}\par\nopagebreak[4]
Manteniendo $n$ y $\sigma$, al aumentar la confianza aumenta el cuantil bilateral $z_{1-\alpha/2}$. Por la fórmula $E=z\sigma/\sqrt n$, \textbf{el error máximo aumenta}: el intervalo se hace más ancho. Esta comparación fija el tamaño muestral; si también cambiara, habría que considerar ambos efectos.
